% Base: Semana 04 Clase 01 y sistemas homogéneos de Semana 05 Clase 01 (2025-A).
\documentclass[compacto,10pt]{pubciencias-notas}

% -- Paquetes adicionales
\usepackage{comandosEPN}

% -- Datos del libro
\autor{Departamento de Formación Básica}
\materia{Álgebra Lineal}
\nota{Semana 03 Clase 01}
\fecha{Semestre 2026-B}

\logouno[1.1cm]{Logo_EPN.eps}
\logodos[1.7cm]{Logo_DFB.eps}

\definecolor{colordef}{cmyk}{0.81,0.62,0.00,0.22}

\begin{document}

\encabezado

\vspace*{-8mm}

\setcounter{section}{1}
\section{Sistemas de ecuaciones}
\setcounter{subsection}{1}
\subsection{Resolución de sistemas lineales}

\setcounter{defi}{2}
\begin{defi}[Sistema consistente e inconsistente]
    Sean $A\in \Mat[\K]{m}{n}$ y $b\in \Mat[\K]{m}{1}$, dado el sistema de ecuaciones lineales
    \[
        Ax=b,
    \]
    se dice que
    \begin{itemize}
    \item
        el sistema es \textbf{inconsistente} si no tiene solución;
    \item
        el sistema es \textbf{consistente} si tiene al menos una solución.
    \end{itemize}
\end{defi}
\begin{teo}
    Un sistema de ecuaciones lineales consistente sobre $\R$ tiene
    una única solución o infinitas soluciones.
\end{teo}

\begin{ejem}\hspace{0pt} 
El sistema de ecuaciones lineales
\[
    \systeme[x_1 x_2]{x_1 -x_2 = 3,
        2x_1 -2x_2 = 4}
\]
no tiene solución, es decir, es un sistema inconsistente. Analicemos la matriz ampliada del sistema
\begin{align*}
    \begin{pmatrix}
    1&-1&|&3\\
    2&-2&|&4
    \end{pmatrix}
    &\sim
    \begin{pmatrix}
    1&-1&|&3\\
    0&0&|&-2
    \end{pmatrix} 
    && -2F_{1} + F_{2} \to F_{2}
\end{align*}
Notemos que la segunda fila de la matriz equivalente implica que
\[
    0x_{1} + 0x_{2} = -2
\]
lo cual es falso, lo que nos permite concluir que el sistema no posee solución, es decir, es inconsistente para cualquier valor de $x_{1}$ y $x_{2}$.
\end{ejem}

\begin{ejem}\hspace{0pt} 
El sistema de ecuaciones lineales
\[
    \systeme[x_1 x_2]{x_1 -x_2 = 3,
        2x_1 -2x_2 = 6}
\]
tiene al menos una solución, es decir, es un sistema consistente. Analicemos la matriz ampliada del sistema
\begin{align*}
    \begin{pmatrix}
    1&-1&|&3\\
    2&-2&|&6
    \end{pmatrix}
    &\sim
    \begin{pmatrix}
    1&-1&|&3\\
    0&0&|&0
    \end{pmatrix} 
    && -2F_{1} + F_{2} \to F_{2}
\end{align*}
La segunda fila corresponde a $0=0$, por lo que el sistema se reduce
a una sola ecuación, $x_1-x_2=3$, con dos variables. Una de ellas puede
tomar cualquier valor real; por tanto,
existirán infinitas soluciones. Más adelante veremos cómo escribirlas
utilizando un parámetro.
\end{ejem}

\begin{ejer}\hspace{0pt} 
Dado el sistema de ecuaciones lineales
\[
    \systeme[x_1 x_2]{2x_1 -x_2 = 3,
        x_1 +x_2 = 0.}
\]
Muestre que es un sistema consistente.
\end{ejer}

\begin{teo}[Teorema de Rouché–Frobenius]\label{sem03cl01:rouche-frobenius}
    Sean $A\in \Mat[\K]{m}{n}$ y $b\in \Mat[\K]{m}{1}$, dado el sistema de ecuaciones lineales
    \[
        Ax=b,
    \]
    se tiene que
    \begin{itemize}
    \item
        el sistema es consistente si y solo si $\rang(A)=\rang(A|b)$;
    \item
        en caso de que el sistema sea consistente, la solución es única si y solo si $\rang(A)=n$; si $\rang(A)<n$, el sistema tiene infinitas soluciones (sobre $\R$).
    \end{itemize}
\end{teo}

\begin{ejem}[Infinitas soluciones: elección de la variable libre]
Resolvamos el sistema de dos ecuaciones con tres incógnitas
\[
    \left\{\begin{aligned}
        x_1+x_2+x_3&=2,\\
        x_2-x_3&=1.
    \end{aligned}\right.
\]
La matriz ampliada ya está escalonada:
\[
    (A\mid b)=\left(\begin{array}{ccc|c}
        1&1&1&2\\0&1&-1&1
    \end{array}\right).
\]
Dado que $n=3$ y $\rang(A)=\rang(A\mid b)=2<n$: el sistema es consistente
con infinitas soluciones y tiene $3-2=1$ variable libre.
\begin{enumerate}[label=\textbf{\arabic*)}]
\item \textbf{Identificar la variable libre.} Dada una matriz escalonada,
los pivotes (primeras entradas no nulas de cada fila) están en las columnas 1 y 2, por lo que $x_1$ y $x_2$
se suelen denominar variables principales. Para aprovechar la forma escalonada, tomamos
$x_3$ como variable libre, pues su columna (columna 3) no tiene pivote, y escribimos
$x_3=t$, con $t\in\R$. Esta elección permite despejar directamente las
variables principales.
\item \textbf{Expresar las demás variables en función de $t$.}
De la segunda ecuación, $x_2=1+x_3=1+t$. Sustituyendo en la primera,
\[
    x_1=2-x_2-x_3=2-(1+t)-t=1-2t.
\]
\item \textbf{Escribir todas las soluciones.}
\[
    x=\begin{pmatrix}1-2t\\1+t\\t\end{pmatrix}
     =\begin{pmatrix}1\\1\\0\end{pmatrix}
      +t\begin{pmatrix}-2\\1\\1\end{pmatrix},\qquad t\in\R.
\]
\textbf{Complete y compruebe:} ¿qué soluciones se obtienen para $t=0$
y $t=1$? Sustitúyalas en las dos ecuaciones originales.
\end{enumerate}
La variable libre puede tomar cualquier valor real; una vez elegido,
las variables principales quedan determinadas.
\end{ejem}

% Material adaptado de Semana 05 Clase 01 (2025-A).
\subsection{Sistemas homogéneos}

\begin{defi}[Sistema homogéneo]
    Sea $A\in\Mat[\R]{m}{n}$. Un sistema de ecuaciones lineales se llama
    \textbf{homogéneo} si todos sus términos independientes son cero,
    es decir, si tiene la forma
    \[
        Ax=0.
    \]
    La solución $x=0$ se llama \textbf{solución trivial}.
    Una solución $x\ne0$, es decir, con al menos una componente distinta
    de cero, se llama \textbf{solución no trivial}.
\end{defi}

Todo sistema homogéneo es consistente, pues $A0=0$. Además, añadir una
columna de ceros no cambia el rango, por lo que $\rang(A)=\rang(A\mid0)$.
\setcounter{prop}{3}
\begin{teo}[Soluciones de un sistema homogéneo]\label{sem03cl01:no-trivial}
Sea $A\in\Mat[\R]{m}{n}$. Para el sistema homogéneo $Ax=0$:
\begin{itemize}
\item Si $\rang(A)=n$, la solución trivial es la única solución.
\item Si $\rang(A)<n$, hay $n-\rang(A)$ variables libres y el sistema
      tiene infinitas soluciones, entre ellas soluciones no triviales.
\end{itemize}
\end{teo}

\begin{ejem}[Soluciones no triviales]
Considere el sistema homogéneo
\[
    \left\{\begin{aligned}
        x_1-x_2&=0,\\
        2x_1-2x_2&=0.
    \end{aligned}\right.
\]
Su matriz ampliada se reduce a
\[
    \left(\begin{array}{cc|c}1&-1&0\\2&-2&0\end{array}\right)
    \sim
    \left(\begin{array}{cc|c}1&-1&0\\0&0&0\end{array}\right),
    \qquad F_2\leftarrow F_2-2F_1.
\]
Como $\rang(A)=1<2$, el sistema tiene soluciones no triviales.
Tomando $x_2=t$, todas sus soluciones son
\[
    x=\begin{pmatrix}t\\t\end{pmatrix},\qquad t\in\R.
\]
Para $t=0$ se obtiene la solución trivial; para $t\ne0$, una solución
no trivial. Por ejemplo, $x=(1,1)^{\mathsf T}$ es una solución no trivial.
\end{ejem}

\begin{ejer}[Clasificación de sistemas homogéneos]
Para cada sistema, determine el rango de la matriz de coeficientes,
indique si existen soluciones no triviales y encuentre todas las soluciones.
\[
\text{a)}\quad
\left\{\begin{aligned}
    2x_1+x_2&=0,\\
    x_1+x_2&=0;
\end{aligned}\right.
\qquad
\text{b)}\quad
\left\{\begin{aligned}
    x_1+2x_2-x_3&=0,\\
    2x_1+4x_2-2x_3&=0.
\end{aligned}\right.
\]
\end{ejer}

\end{document}
